\documentclass[11pt,a4paper]{article}
\usepackage[T1]{fontenc}
\usepackage[utf8]{inputenc}
\usepackage{amsmath,amssymb,amsthm}
\usepackage[margin=1in]{geometry}
\usepackage[colorlinks=true,linkcolor=blue,urlcolor=blue]{hyperref}
\theoremstyle{plain}
\newtheorem{theorem}{Theorem}
\newtheorem{lemma}[theorem]{Lemma}
\newtheorem{proposition}[theorem]{Proposition}
\newtheorem{corollary}[theorem]{Corollary}
\theoremstyle{definition}
\newtheorem{remark}[theorem]{Remark}

\title{The empirical closed form for OEIS A199834, proved from an exact model}
\author{Adrian Perez Fontelles\\ \small Independent researcher}
\date{13 September 2026}
\begin{document}
\maketitle

\begin{abstract}
OEIS A199834 carries an empirical closed form: a polynomial in $n$ of degree $4$, recorded on the
entry and never marked settled. It is true. The entry's count is given exactly by an Ehrhart
quasi-polynomial, from which $a$ provably satisfies a monic linear recurrence of order $S=6$
derived below; a polynomial of degree $4$ satisfies $(z-1)^{4+1}$; so the difference
satisfies the product, and $6+4+1$ consecutive agreements settle it. No transfer matrix is
involved and the count is not a walk.
\end{abstract}

\noindent\small 2020 Mathematics Subject Classification. 05A15, 05A19, 11P21.\normalsize

\section{The conjecture}

OEIS A199834 is ``Number of -n..n arrays x(0..4) of 5 elements with zero sum and no two neighbors summing to zero.''. It has offset $1$ and begins
\[
4, 114, 646, 2146, 5390, 11384, 21364, 36796,\ \dots
\]
The entry states, as a conjecture contributed by its author and never marked settled:
\begin{quote}
Empirical: a(n) = (115/12)*n\^{}4 - (29/6)*n\^{}3 + (5/12)*n\^{}2 - (7/6)*n.
\end{quote}
As of the ``Last modified'' line on the live entry (Oct 03 2025 08:20:04, revision 11)
this is still recorded as empirical, and nothing on the entry records it as proved.

\section{The count is a quasi-polynomial}

The array has a FIXED number $L$ of entries and an alphabet that grows with $n$. Every
condition the entry imposes is a boolean combination of statements
$\sum_i c_ix_i + c\,n \bowtie 0$ with integer coefficients and $\bowtie$ one of
$=,\ne,<,\le,>,\ge$: the bounds $|x_i|\le n$ or $0\le x_i\le n$, a vanishing sum, a difference
that must not vanish, a pair that must total exactly $n$. Each is homogeneous in $(x,n)$
jointly, so the admissible $x$ for a given $n$ are precisely the integer points at height $n$
of a finite union of relatively open rational cones in $\mathbb{R}^{L+1}$, and $a(n)$ is
the Ehrhart quasi-polynomial of that union: of degree at most $L$, with period dividing the
heights of the cones' ray generators.

\section{The bound}

The period is derived, not assumed. Every ray of every cell of the arrangement is cut out by
$L$ linearly independent hyperplanes taken from the constraint set; writing that $L\times(L+1)$
system as $M$, its null direction has $t$-component $\pm\det$ of the $x$-parts, so the
primitive generator $(x^*,t^*)$ has $t^*$ dividing that determinant. Let $T$ be the set of
those determinants. A simplicial cone in a triangulation has at most $L+1$ rays, each
contributing a factor $1-z^{t_i}$ with $t_i\in T$, so
\[
A(z)\;=\;\prod_{d\,\mid\,t\ \text{for some}\ t\in T}\Phi_d(z)^{L+1}
\]
annihilates $a$, and $S=\deg A=6$. A congruence condition modulo $m$ is counted one residue
class at a time; a class is a coset of $m\mathbb{Z}^L$, and a generator primitive in
$\mathbb{Z}^{L+1}$ must be scaled by a divisor of $m$ to lie in it, so each $t$ is replaced
by $mt$ and nothing else changes.

Over the common denominator the numerator has degree below that of $A$ for every cell of the
arrangement that HAS a ray, because such a cell's numerator collects a fundamental
parallelepiped whose heights are below the sum of its generators' heights. The origin is a cell
too and it has none: its series is the constant $1$, which over $A$ is $A/A$ and reaches degree
exactly $\deg A$. The bound in force is therefore $zA(z)$, of order $S=6$, and
$a(0),\dots,a(S-1)$, computed exactly, determine every later term. Nothing is fitted, and the
annihilator is tested on terms the model was not asked for before anything is claimed --- a
bound one too small fails that test, which is how this one was found.

\section{The proof}

\begin{theorem}
The empirical closed form holds for every $n$.
\end{theorem}

\begin{proof}
Write $P$ for the claimed polynomial and $r(n)=a(n)-P(n)$. The annihilator $A$ derived above
kills $a$, and $(z-1)^{4+1}$ kills $P$, so $A(z)(z-1)^{4+1}$ --- monic, of order
$6+4+1$ --- kills $r$. The model was evaluated exactly in integer arithmetic and $P$ in
exact rational arithmetic at every index from $n=1$ to $n=1+6+4+1$, and $r$ vanished
at all of them. A monic recurrence with $6+4+1$ consecutive zeros has every later term zero,
so $r\equiv0$.
\end{proof}

\section{Verification}

Three checks, each independent of the algebra above.

First, the model reproduces all $31$ terms the entry publishes, exactly, in integer
arithmetic. This is what ties the model to the entry: it is built by reading the entry's
English, and a misreading gives different counts.

Second, the claimed polynomial was evaluated directly on those published terms, in exact
rational arithmetic and with no model involved, and agrees with every one of them.

Third, the derived annihilator was itself tested on terms the model was not asked for: the
model was evaluated beyond the $S$ values that determine it and $A$ reproduced them. A bound
that is too small fails this test, which is why it is run.

\begin{thebibliography}{9}
\bibitem{oeis} The OEIS Foundation, \emph{The On-Line Encyclopedia of Integer
Sequences}, \url{https://oeis.org}, 2026. Sequence A199834.
\bibitem{beckrobins} M.~Beck and S.~Robins, \emph{Computing the Continuous Discretely},
2nd ed., Springer, 2015. (Ehrhart quasi-polynomials and their periods.)
\bibitem{stanley} R.~P.~Stanley, \emph{Enumerative Combinatorics}, Volume 1, 2nd ed.,
Cambridge University Press, 2011.
\end{thebibliography}

\end{document}
